% titulo: Referencias
% estado: borrador
% actualizado: 2026-10-19
% Formato IEEE. Cada entrada se cita como \cite{refN}; la numeración la asigna LaTeX en este orden.
% Para agregar una referencia: añade un \bibitem al final y cítalo en el texto.

\begin{thebibliography}{99}
\setlength{\itemsep}{4pt}
  \bibitem{ref1} S. Moutinho, “Latin American journals are open-access pioneers. Now, they need an audience,” Science, vol. 386, no. 6726, Dec. 5, 2024, doi: 10.1126/science.zjn80sl. Disponible en: \url{https://www.science.org/content/article/latin-american-journals-are-open-access-pioneers-now-they-need-audience}

  \bibitem{ref2} A. Casas Niño de Rivera, ``Opportunities for Scholarly Publishing Services in Latin America,'' M.Pub. thesis, Faculty of Communication, Art and Technology, Simon Fraser University, Burnaby, BC, Canada, Fall 2023. Disponible en: \url{https://summit.sfu.ca/_flysystem/fedora/2024-02/etd22832.pdf}

  \bibitem{ref3} M. Khalifa and M. Albadawy, “Using artificial intelligence in academic writing and research: An essential productivity tool,” Computer Methods and Programs in Biomedicine Update, vol. 5, p. 100145, Jan. 2024, doi: 10.1016/j.cmpbup.2024.100145. Disponible en: \url{https://www.sciencedirect.com/science/article/pii/S2666990024000120}

  \bibitem{ref4} P. Accuosto, F. Ronzano, D. Ferrés y H. Saggion, “Multi-level mining and visualization of scientific text collections: Exploring a bi-lingual scientific repository,” en Proceedings of the 6th International Workshop on Mining Scientific Publications (WOSP 2017), Toronto, Canada, 2017, pp. 9--16, doi: 10.1145/3127526.3127529. Disponible en: \url{https://dl.acm.org/doi/10.1145/3127526.3127529}

  \bibitem{ref5} K. A. Yaghi, “Using artificial intelligence for formatting academic journal papers,” American Academic \& Scholarly Research Journal, vol. 16, no. 4, pp. 28--57, Oct. 2024. Disponible en: \url{https://aasrc.org/aasrj/index.php/aasrj/article/view/2296}

  \bibitem{ref6} L. W. Markewich, “Extracting structured information from academic documents,” M.Sc. thesis, University of Saskatchewan, Saskatoon, SK, Canada, Apr. 11, 2022. https://harvest.usask.ca/items/d46f9406-8b3c-4999-9b39-12f472e06c4a

  \bibitem{ref7} S. Nigam, S. Verma, and P. Nagabhushan, “Document Analysis and Recognition: A Survey,” TechRxiv, Mar. 27, 2023, doi: 10.36227/techrxiv.22336435.v1. Disponible en: \url{https://www.techrxiv.org/doi/10.36227/techrxiv.22336435}

  \bibitem{ref8} DM. Minouei, M. R. Soheili, and D. Stricker, “Document layout analysis with an enhanced object detector,” in 2021 5th International Conference on Pattern Recognition and Image Analysis (IPRIA), Apr. 2021, pp. 1--5, doi: 10.1109/IPRIA53572.2021.9483509. Disponible en: \url{https://ieeexplore.ieee.org/document/9483509}

  \bibitem{ref9} S. M. B. W. Dissanayake and P. P. G. D. Asanka, “A stage-aware framework for automated IEEE research paper compliance checking using layout analysis, rule-based validation, and semantic evaluation,” in 2026 6th International Conference on Advanced Research in Computing (ICARC), Feb. 2026, pp. 1--6, doi: 10.1109/ICARC68737.2026.11454060. Disponible en: \url{https://www.researchgate.net/publication/403371900_A_Stage-Aware_Framework_for_Automated_IEEE_Research_Paper_Compliance_Checking_Using_Layout_Analysis_Rule-Based_Validation_and_Semantic_Evaluation}

  \bibitem{ref10} V. D. Tran, M. L. Nguyen, K. Shirai, and K. Satoh, “An approach of rhetorical status recognition for judgments in court documents using deep learning models,” in 2019 11th International Conference on Knowledge and Systems Engineering (KSE), Oct. 2019, pp. 1--6, doi: 10.1109/KSE.2019.8919370. Disponible en: \url{https://ieeexplore.ieee.org/document/8919370}

  \bibitem{ref11} S. Dündar, “Performance analysis of text chunking methods for LLM-based document summarization,” IEEE Xplore, Sep. 2025. Disponible en: \url{https://ieeexplore.ieee.org/document/11206896}

  \bibitem{ref12} J. Ye and K. Xu, “UDMtext: A unified deep mutual learning-based framework for text detection in complex documents,” in 2025 6th International Conference on Computer Vision, Image and Deep Learning (CVIDL), May 2025, doi: 10.1109/CVIDL65390.2025.11086049. Disponible en: \url{https://www.researchgate.net/publication/394005168_UDMtext_a_Unified_Deep_Mutual_Learning-Based_Framework_for_Text_Detection_in_Complex_Documents}

  \bibitem{ref13} M. Moskalev and T. T. Gazizov, “AI-based algorithm for significant information extraction from unstructured textual data,” in 2025 International Russian Automation Conference (RusAutoCon), Sep. 2025, doi: 10.1109/RusAutoCon65989.2025.11177408. Disponible en: \url{https://www.researchgate.net/publication/396022201_AI-Based_Algorithm_for_Significant_Information_Extraction_from_Unstructured_Textual_Data}

  \bibitem{ref14} R. Poorvadevi, H. Surendar, and S. SriRamakrishnan, “Effective analysis of data processing on cloud using serverless computing,” in 2025 8th International Conference on Trends in Electronics and Informatics (ICOEI), Apr. 2025, doi: 10.1109/ICOEI65986.2025.11013090. Disponible en: \url{https://www.researchgate.net/publication/392245936_Effective_Analysis_of_Data_Processing_on_Cloud_using_Serverless_Computing}

  \bibitem{ref15} J. Chen, S. Yang, and S. Jiang, “Automatic classification and recognition of complex documents based on Faster RCNN,” in 2019 14th IEEE International Conference on Electronic Measurement \& Instruments (ICEMI), Nov. 2019, pp. 573--577, doi: 10.1109/ICEMI46757.2019.9101847. Disponible en: \url{https://www.researchgate.net/publication/341695915_Automatic_classification_and_recognition_of_complex_documents_based_on_Faster_RCNN}

  \bibitem{ref16} S. Takeshita, T. Green, N. Friedrich, K. Eckert, and S. P. Ponzetto, “X-SCITLDR: Cross-lingual extreme summarization of scholarly documents,” in Proceedings of the ACM/IEEE Joint Conference on Digital Libraries (JCDL 2022), Cologne, Germany, Jun. 20--24, 2022, doi: 10.1145/3529372.3530938.

  \bibitem{ref17} W. Ma, M. Cao, J. Ma, Z. Dong, C. Yang, and Z. Li, “Mamba-YOLO: Multi-level adaptive rectangular convolution for document layout analysis,” Pattern Recognition, vol. 170, Art. no. 112031, Feb. 2026, doi: 10.1016/j.patcog.2025.112031.

  \bibitem{ref18} L. Pisaneschi, A. Gemelli, and S. Marinai, “Automatic generation of scientific papers for data augmentation in document layout analysis,” Pattern Recognition Letters, vol. 167, pp. 38--44, 2023, doi: 10.1016/j.patrec.2023.01.018.

  \bibitem{ref19} X. Wu et al., “Cross-domain document layout analysis using document style guide,” Expert Systems with Applications, vol. 245, Art. no. 123039, Jul. 2024, doi: 10.1016/j.eswa.2023.123039.

  \bibitem{ref20} A. Peña et al., “Continuous document layout analysis: Human-in-the-loop AI-based data curation, database, and evaluation in the domain of public affairs,” Information Fusion, vol. 108, Art. no. 102398, Aug. 2024, doi: 10.1016/j.inffus.2024.102398.

  \bibitem{ref21} P. Mishra, “Domain adaptive learning for document layout analysis and object detection using classifier alignment mechanism,” Signal Processing: Image Communication, vol. 116, Art. no. 116986, Aug. 2023, doi: 10.1016/j.image.2023.116986.

  \bibitem{ref22} N. Raman, S. Shah, and M. Veloso, “Synthetic document generator for annotation-free layout recognition,” Pattern Recognition, vol. 128, Art. no. 108660, Aug. 2022, doi: 10.1016/j.patcog.2022.108660.

  \bibitem{ref23} K. Nandi and S. S. Sathya, “Visual document understanding: A comparative review of modern methods,” in Computer Vision and Image Processing, Communications in Computer and Information Science, pp. 411--427, 2026, doi: 10.1007/978-3-031-93694-4\_29.

  \bibitem{ref24} C. Xu, Y. Li, S. Cao, H. Zhang, H. Bi, and Y. Chen, “HiM: Hierarchical multimodal network for document layout analysis,” Applied Intelligence, vol. 53, no. 20, pp. 24314--24326, Jul. 2023, doi: 10.1007/s10489-023-04782-3.

  \bibitem{ref25} R.-Y. Cao, Y.-X. Cao, G.-B. Zhou, and P. Luo, “Extracting variable-depth logical document hierarchy from long documents,” Journal of Computer Science and Technology, vol. 37, no. 3, pp. 699--718, May 2022, doi: 10.1007/s11390-021-1076-7.

  \bibitem{ref26} P. Constantopoulos and V. Pertsas, “From publications to knowledge graphs,” in Communications in Computer and Information Science, 2020, pp. 18--33, doi: 10.1007/978-3-030-44900-1\_2.

  \bibitem{ref27} Z. Chkirbene, R. Hamila, A. Gouissem, and U. Devrim, “Large language models (LLM) in industry: A survey of applications, challenges, and trends,” in 2024 19th International Conference on High Performance Computing \& Simulation / HONET, Dec. 2024, pp. 229--234, doi: 10.1109/HONET63146.2024.10822885.

  \bibitem{ref28} L. Wang et al., “A survey on large language model based autonomous agents,” arXiv preprint arXiv:2308.11432, Aug. 2023.

  \bibitem{ref29} D. J. Anderson, Kanban: Successful Evolutionary Change for Your Technology Business. Sequim, WA, USA: Blue Hole Press, Apr. 2010.

\end{thebibliography}

% Pendientes de cita (añadir como \bibitem y citar):
% - Guía de Scrum (Schwaber y Sutherland) — 2.1.3 y 3.2
% - Método RICE (Intercom) — 2.1.3 y 3.2
% - Spec-Driven Development — 2.1.3 y 3.2
% - Documentación de OJS 3.x / PKP (API REST y plugins)
% - Ley N.º 29733, Ley de Protección de Datos Personales (Perú) — 3.11
% Observaciones: ref3 (Khalifa y Albadawy) no se cita en el cuerpo; ref29 (Kanban) quedó sin uso al adoptar Scrum.